\documentclass[11pt]{article}

\usepackage[margin=1in]{geometry}
\usepackage{amsmath, amssymb}
\usepackage{booktabs}
\usepackage{longtable}
\usepackage{array}
\usepackage[hidelinks]{hyperref}
\usepackage{xcolor}
\usepackage{enumitem}
\usepackage{microtype}

\newcommand{\code}[1]{\texttt{#1}}
\newcommand{\term}[1]{\textit{#1}}

\title{\textbf{ResonantOS Pi-Native Session Credential}\\[2mm]
\large P2 --- User Live Proof: Fixes and Results}
\author{ResonantOS Engineering}
\date{September 29, 2026}

\begin{document}

\maketitle

\begin{abstract}
This document records the implementation and live validation of the
Pi-native session-credential path (P2) in ResonantOS. P2 proves a single,
reviewed behavior: a provider credential entered \emph{once} through the normal
ResonantOS Settings interface can securely power a real Pi model session,
delivered to the native Pi process through a session-scoped environment
variable, with no durable copy of the credential written into Pi. The live proof
\emph{passed}: a Settings-configured OpenRouter credential powered a real Pi
process that invoked a real model through OpenRouter, returned the expected
response, left Pi's \code{auth.json} byte-identical, and failed closed with
\code{permission-denied} after revocation.
\end{abstract}

\tableofcontents

\section{Objective}

Prepare and run the first user-controlled, end-to-end proof that a real Pi
process can consume an AI-provider credential configured once through
ResonantOS Settings and successfully call a real model, without duplicating or
persisting that credential inside Pi.

The acceptance chain is:

\begin{center}
\begin{tabular}{c}
\term{User enters credential once} $\rightarrow$ \term{ROS Settings}
$\rightarrow$ \term{Provider Profile / Vault} $\rightarrow$
\term{Pi authorization} $\rightarrow$ \term{session-scoped credential}
$\rightarrow$ \term{REAL Pi process} $\rightarrow$
\term{REAL AI provider/model} $\rightarrow$ \term{REAL response}
$\rightarrow$ \term{Pi terminates} $\rightarrow$
\term{session authority cleaned up}
\end{tabular}
\end{center}

\noindent with \emph{no durable Pi credential copy}.

\section{Starting Point}

\begin{itemize}
  \item \textbf{Candidate SHA:} \code{8ab442c84f76298649cf5670356b3ccb55cdd73f}
  \item \textbf{Base branch:} \code{feature/pi-native-session-credential-p1-1-protocol-gate}
  \item \textbf{Work branch:} \code{feature/pi-native-session-credential-p2}
  \item \textbf{Commit:} \code{d545edae} ---
    \code{feat(pi-native): P2 session-credential launcher and host wiring}
\end{itemize}

The already-audited chain (Provider Profile exists $\rightarrow$ host-derived
protocol compatible $\rightarrow$ native Pi provider mapping $\rightarrow$
selected model belongs to profile $\rightarrow$ binding + grant authorization
$\rightarrow$ approved Pi executable $\rightarrow$ credential resolution
$\rightarrow$ private session environment) was preserved and extended, not
redesigned.

\section{Changes Delivered}

\subsection{P2-A --- Host Wiring}

A new module \code{pi-native-session-service.mjs} composes the plan-only
credential adapter (\code{pi-native-credential-adapter.mjs}) with the bounded
launcher so the private launch plan flows directly:

\begin{center}
\code{host} $\rightarrow$ \code{Pi credential planner} $\rightarrow$
\code{private launch material} $\rightarrow$ \code{process launcher}
\end{center}

The host supplies every authority and nothing is read from a manifest or caller:
authoritative provider profiles and model catalog (provider host), credential
resolution from the ROS provider/vault mechanism, registry binding and grant
authorization, the reviewed \code{piCommand()} executable, and a host-approved
project path. The raw, secret-bearing plan never leaves the host process; only
\code{redactLaunchPlan()} output crosses an observability/UI boundary.

\subsection{P2-B --- Bounded Process Launcher}

A new module \code{pi-process-launcher.mjs} implements the smallest bounded
launcher required for the proof. It:

\begin{itemize}
  \item spawns only the executable returned by \code{piCommand()};
  \item sets \code{cwd} to the host-approved project path;
  \item builds argv as exactly
    \code{[--print, --no-session, \ldots{}provider/model args, prompt]};
  \item never uses \code{--api-key};
  \item uses the private session-environment plan as the child environment
    (no blanket \code{process.env} inheritance);
  \item enforces a bounded timeout with deterministic \code{SIGTERM}
    $\rightarrow$ \code{SIGKILL} cleanup;
  \item bounds and sanitizes stdout/stderr, stripping the credential from any
    captured text;
  \item disables shell interpolation and rejects arbitrary manifest commands or
    caller argv.
\end{itemize}

Pi's non-interactive headless mechanism (\code{--print}) is used for the proof;
ephemeral \code{--no-session} prevents durable session writes.

\subsection{Executable Allowlist Extension}

\code{pi-runtime.mjs} was extended to resolve \code{nvm}-managed install roots
(\code{\textasciitilde/.nvm/versions/node/<version>/\{bin,lib/node\_modules\}}),
mirroring the existing npm-global discipline, so the real Pi installation is
resolved by the reviewed allowlist.

\subsection{Proof Trigger}

\code{run-bridge-minimal.mjs} instantiates the service with real authorities and
adds a loopback-only, bridge-token and capability-gated route
\code{POST /pi-native/proof} that runs the proof in-process (sharing the
Settings-configured session credential) and returns only redacted evidence.

\section{Environment Findings}

The live environment differed from the task assumptions in ways that were
recorded and worked around, but that did not block the proof:

\begin{itemize}
  \item \textbf{Pi version skew.} The installed binary reports \code{0.74.2};
    the user's Pi config records \code{lastChangelogVersion: 0.80.3}. The six
    provider environment variables in the host mapping are still present in
    \code{0.74.2}.
  \item \textbf{Install location.} Pi is installed under
    \code{nvm} (\code{node/v22.13.0}), not an npm-global or fixed system root;
    hence the allowlist extension above.
  \item \textbf{Built-in providers.} \code{0.74.2} exposes
    \code{minimax}, \code{ollama}, \code{openrouter}, \code{xai}, and
    \code{zai} as built-in providers; \code{openai} and \code{deepseek} are not
    built-in in this version.
  \item \textbf{Model catalog mismatch.} The ROS catalog lists
    \code{openai/gpt-5.5} for OpenRouter, but \code{0.74.2} does not include
    that model. A model present in both the ROS catalog and Pi's catalog was
    selected instead (\code{anthropic/claude-sonnet-4.5}).
  \item \textbf{Local-model fallback risk.} Pi's default provider is a local
    \code{ollama} model. A discriminating identity prompt during the proof
    confirmed the response originated from an Anthropic Claude model via
    OpenRouter, not the local fallback.
\end{itemize}

\section{Live Proof Results}

\begin{longtable}{@{}l p{9.5cm}@{}}
\toprule
\textbf{Field} & \textbf{Value} \\
\midrule
\endhead
Outcome & \textbf{PASS} \\
Candidate SHA & \code{8ab442c84f76298649cf5670356b3ccb55cdd73f} \\
Pi version & \code{0.74.2} \\
Provider Profile & \code{openai-compatible-ros-openrouter-test-api} (label
  \code{ros-openrouter-test-api}) \\
Provider identity & OpenRouter (\code{templateId: openrouter}) \\
Derived protocol & \code{openai-compatible} \\
Model & \code{anthropic/claude-sonnet-4.5} \\
Credential configured in ROS Settings & yes \\
Credential mechanism & \code{session-environment} \\
Real Pi process launched & yes (exit \code{0}, $\approx$22\,s) \\
Real inference performed & yes \\
Expected response received & yes --- \code{ROS\_PI\_NATIVE\_CREDENTIAL\_PROOF\_OK} \\
Environment variable (name only) & \code{OPENROUTER\_API\_KEY} \\
\bottomrule
\end{longtable}

\section{Credential Hygiene and Leakage Verification}

\begin{longtable}{@{}l l@{}}
\toprule
\textbf{Check} & \textbf{Result} \\
\midrule
\endhead
Secret in argv & no \\
Secret in manifest & no \\
Secret in Git & no \\
Secret in public projection & no (env key \emph{names} only) \\
Secret in logs / evidence & no \\
Secret written to \code{auth.json} & no \\
Credential duplicated into Pi & no \\
\bottomrule
\end{longtable}

Pi's \code{auth.json} was compared before and after the live proof and remained
unchanged:

\begin{center}
\begin{tabular}{@{}l l@{}}
\toprule
\textbf{Property} & \textbf{Value} \\
\midrule
Size & 557 bytes \\
mtime & 1783522429 (2026-07-08 10:53:49) \\
SHA-256 & \code{ac488b7304fa0fba1590769c552f4802b9a23753581267692a0405d38fc73ac0} \\
\bottomrule
\end{tabular}
\end{center}

\noindent The credential flowed only through: Settings $\rightarrow$ session
memory $\rightarrow$ session environment (\code{OPENROUTER\_API\_KEY})
$\rightarrow$ the Pi process. It was never duplicated or persisted.

\section{Revocation Proof}

After the successful model turn, Pi's \code{agent-runtime} grant was revoked
through the host registry. A second native Pi launch plan was then attempted:

\begin{center}
\code{\{"ok":false,"code":"permission-denied"\}} \quad (HTTP 403)
\end{center}

The underlying provider profile and credential were \emph{not} revoked or
deleted; the shared provider credential remained \code{configured: true} and
usable by other authorized ROS consumers.

\section{Cleanup}

\begin{itemize}
  \item Pi process terminated (no lingering \code{pi} processes).
  \item Session environment was in-memory only.
  \item Pi's \code{agent-runtime} grant was revoked.
  \item Provider profile and credential left in place.
\end{itemize}

\section{Files Changed}

\subsection{Modified}

\begin{itemize}
  \item \code{browser-first/host/pi-runtime.mjs} (nvm install-root support)
  \item \code{browser-first/host/run-bridge-minimal.mjs} (host wiring and
    \code{/pi-native/proof} route)
  \item \code{browser-first/test/pi-native-session-credential.test.mjs}
  \item \code{docs/architecture/MODULE-OWNERSHIP.md}
\end{itemize}

\subsection{Added}

\begin{itemize}
  \item \code{browser-first/host/pi-process-launcher.mjs}
  \item \code{browser-first/host/pi-native-session-service.mjs}
  \item \code{browser-first/test/pi-process-launcher.test.mjs}
  \item \code{browser-first/test/pi-native-session-service.test.mjs}
\end{itemize}

\section{Tests}

All deterministic checks pass (Node 24.21.0, \code{node --test}):

\begin{itemize}
  \item 62 tests across the harness, provider, Pi credential, launcher, and
    service suites.
  \item \code{bridge-route-capability-audit} (route/capability boundary).
  \item \code{bridge-capability-token-consistency} (capability catalog drift
    guard).
  \item \code{module-ownership-doc} (ownership map coverage).
\end{itemize}

\section{Known Limitations}

\begin{enumerate}
  \item \textbf{Pi version skew.} The installed binary (\code{0.74.2}) differs
    from the configured/expected version (\code{0.80.3}). Re-validate the
    provider mapping against the intended Pi version when it is installed.
  \item \textbf{Provider availability.} \code{openai} and \code{deepseek} are
    not built-in providers in \code{0.74.2}.
  \item \textbf{Model catalog drift.} The ROS catalog model
    \code{openai/gpt-5.5} is not present in Pi \code{0.74.2}; a mutually
    available model was used for the proof.
  \item \textbf{Authorization conflation.} The P2 authorization gate authorizes
    the native launch through the harness \code{primary-agent} slot, which
    temporarily displaced the Augmentor chat in the isolated test registry. A
    cleaner follow-up authorizes on installed + enabled + granted + binding via
    \code{registry.snapshot()} without requiring slot ownership.
  \item \textbf{Model self-identification.} The requested model self-identified
    with an earlier version label when asked its name (a common LLM quirk); the
    response was nonetheless confirmed to originate from a real Anthropic model
    via OpenRouter, not a local fallback.
  \item \textbf{UI not yet implemented.} The session-spin and Pi add-on
    connection UI are intentionally out of scope for P2 and deferred to a later
    task.
\end{enumerate}

\section{Conclusion}

P2 is complete and \emph{passes}. A credential entered once in ResonantOS
Settings powered a real Pi model session through a session-scoped
\code{OPENROUTER\_API\_KEY} environment variable, with no durable Pi credential
copy, no secret in any projection or log, an unchanged \code{auth.json}, and a
\code{permission-denied} revocation. The exact product behavior required by the
acceptance criterion was demonstrated.

\end{document}
